%%%% ijcai26.tex

\typeout{IJCAI--ECAI 26 Instructions for Authors}

% These are the instructions for authors for IJCAI--ECAI 26.

\documentclass{article}
\pdfpagewidth=8.5in
\pdfpageheight=11in

% The file ijcai26.sty is a copy from ijcai22.sty
% The file ijcai22.sty is NOT the same as previous years'
\usepackage{ijcai26}

% Use the postscript times font!
\usepackage{times}
\usepackage{soul}
\usepackage{url}
\usepackage[hidelinks]{hyperref}
\usepackage[utf8]{inputenc}
\usepackage[small]{caption}
\usepackage{graphicx}
\usepackage{amsmath}
\usepackage{amsthm}
\usepackage{booktabs}

\usepackage{algorithm}
\usepackage{algpseudocode}
\usepackage[switch]{lineno}



%%%%%%%%%%%%%


% \usepackage{microtype}
% \usepackage{algorithm}
 % \usepackage[ruled,linesnumbered]{algorithm2e}
% \usepackage{algpseudocode}
% \usepackage{natbib}
\usepackage{varwidth}
\usepackage{booktabs}
\usepackage{multirow}
\usepackage{amsmath}
\usepackage{rotating}
\usepackage{amssymb}
\usepackage{subcaption}
\usepackage{array}
\usepackage{tabularx}
\usepackage{ragged2e}
% \raggedbottom
\usepackage{amsfonts}
\usepackage{newtxmath}
\usepackage{dsfont}
\usepackage{tcolorbox}
\usepackage{enumitem}
% \usepackage{algpseudocode}

\usepackage{tcolorbox}
\usepackage{listings}

\renewcommand{\algorithmicrequire}{\textbf{Input:}}
\renewcommand{\algorithmicensure}{\textbf{Output:}}

%%%%%%%%

% Comment out this line in the camera-ready submission
\linenumbers

\urlstyle{same}

% the following package is optional:
%\usepackage{latexsym}

% See https://www.overleaf.com/learn/latex/theorems_and_proofs
% for a nice explanation of how to define new theorems, but keep
% in mind that the amsthm package is already included in this
% template and that you must *not* alter the styling.
\newtheorem{example}{Example}
\newtheorem{theorem}{Theorem}

% Following comment is from ijcai97-submit.tex:
% The preparation of these files was supported by Schlumberger Palo Alto
% Research, AT\&T Bell Laboratories, and Morgan Kaufmann Publishers.
% Shirley Jowell, of Morgan Kaufmann Publishers, and Peter F.
% Patel-Schneider, of AT\&T Bell Laboratories collaborated on their
% preparation.

% These instructions can be modified and used in other conferences as long
% as credit to the authors and supporting agencies is retained, this notice
% is not changed, and further modification or reuse is not restricted.
% Neither Shirley Jowell nor Peter F. Patel-Schneider can be listed as
% contacts for providing assistance without their prior permission.

% To use for other conferences, change references to files and the
% conference appropriate and use other authors, contacts, publishers, and
% organizations.
% Also change the deadline and address for returning papers and the length and
% page charge instructions.
% Put where the files are available in the appropriate places.


% PDF Info Is REQUIRED.

% Please leave this \pdfinfo block untouched both for the submission and
% Camera Ready Copy. Do not include Title and Author information in the pdfinfo section
\pdfinfo{
/TemplateVersion (IJCAI.2026.0)
}

\title{MLLMs Get It Right, Then Get It Wrong: Tracing and Correcting \\ Late-Layer Textual Bias}

% Single author syntax
\author{
    Author Name
    \affiliations
    Affiliation
    \emails
    email@example.com
}

% Multiple author syntax (remove the single-author syntax above and the \iffalse ... \fi here)
\iffalse
\author{
First Author$^1$
\and
Second Author$^2$\and
Third Author$^{2,3}$\And
Fourth Author$^4$\\
\affiliations
$^1$First Affiliation\\
$^2$Second Affiliation\\
$^3$Third Affiliation\\
$^4$Fourth Affiliation\\
\emails
\{first, second\}@example.com,
third@other.example.com,
fourth@example.com
}
\fi

% Define a switch to show/hide revisions
\newif\ifshowrevisions
\showrevisionstrue  % Change to \showrevisionsfalse to hide blue color


% Define the revision environment
\newenvironment{revised}{%
  \ifshowrevisions%
    % \color{blue}%
  \fi%
}{}

\begin{document}

\maketitle

% \begin{abstract}
%     \input{chapters/0-abstract}
% \end{abstract}

% % \section{Introduction}

% 

Multimodal AI systems are moving into real-world use. They help diagnose diseases, moderate content, and assist drivers. This growing deployment raises a pointed question: \emph{can we trust what they see?} Imagine a radiologist reviewing a chest X-ray with an AI assistant. The system receives both the image and the radiologist's preliminary notes. If those notes contain an error, will the AI catch it by looking at the image? Or will it simply echo the mistake?

Unfortunately, the evidence points to the second outcome. Show a model an image of a red car alongside text claiming ``the car is blue,'' and most MLLMs will confidently answer ``blue''~\cite{Insight_Over_Sight,PhD_ChatGPT,wu2024autohallusion}. They
trust the text over their own visual perception. This pattern appears
across different model families~\cite{llava-1.5,instructblip2023,bai2025qwen2,bai2025qwen3vltechnicalreport,chen2024internvl} and benchmarks~\cite{zhang2025robust,jia2025benchmarking,qian2024easy}. In safety-critical applications~\cite{bannur2024maira,hou2025one}, such blind trust in text could cause serious harm.

\begin{figure}[t]
  \centering
  \includegraphics[width=\linewidth]{figures/teaser.pdf}
  \caption{How predictions shift under visual-textual conflict. (a) The model sees a red car but outputs the text-suggested answer. (b) Modal Dominance Ratio (MDR) across layers: early layers favors vision, but late layers override toward text. (c) Shift direction predicts correctness—text-ward shifts correlate with failures.}
  \label{fig:teaser}
\end{figure}

\emph{Why} does this happen? Prior work describes the \textit{\underline{behavior}} but not the \textit{\underline{mechanism}}~\cite{Insight_Over_Sight,zhang2025robust}. Is visual information poorly encoded in the first place? Is it lost during cross-modal fusion? Or is it correctly processed but later overwritten? Understanding the failure mode matters: different causes require different solutions.


To find out, we trace how a model's predictions evolve across its layers. Following~\cite{depth_adaptive2020,confident_adaptive2022}, we project hidden states at each layer to output probabilities. This lets us observe what the model ``prefers" at each depth, and what we find is surprising. Surprisingly, in many conflict failures, \textbf{the model assigns higher probability to the correct visual answer in intermediate layers}. It then reverses course and outputs the text answer (Figure~\ref{fig:teaser}(a-b)).


We term this phenomenon \textbf{late-layer textual override}. To measure it, we introduce
Modal Dominance Ratio (MDR), which quantifies vision-versus-text
preference at each layer. In failure cases, MDR starts positive (visual)
and flips negative (textual) in later layers. The visual information was there; it simply didn't survive the full forward pass.

\textbf{Not every late-layer shift is harmful.}
Such prediction shift occurs in successful cases too~\cite{wang2025shift}, not just failures. What distinguishes the two outcomes is the \emph{direction} of the shift. We quantify this using the change in MDR across the transition point. As shown in Figure~\ref{fig:teaser}(c), in failures, 85\% of transitions move toward text. In successes, 89\% move toward vision. Late layers sometimes help; they sometimes hurt. The key
is knowing which.

This directional asymmetry points toward a practical solution. If a model holds a confident visual prediction at the transition layer but abandons it by the final layer, we have reason to suspect harmful override. We capture this intuition with two signals derived from our analysis. \textit{Anchor confidence} measures how certain the transition-layer prediction is. \textit{Prediction retention} captures whether it survives to the output. A high anchor confidence paired with low retention is the signature of harmful override.

\textbf{Recovering what was known.} These observations motivate \textbf{\underline{C}}onflict-\textbf{\underline{A}}ware \textbf{\underline{L}}ayer \textbf{\underline{R}}eference \textbf{\underline{D}}ecoding (CALRD), a training-free method that restores intermediate predictions when override is detected. Our core insight is that predictions at the transition layer still carry the vision-grounded answer before it gets overwritten. By blending transition-layer logits back into the final distribution, we can recover it. CALRD first locates the layer where the output distribution shifts most sharply. It then uses anchor confidence and prediction retention to set correction strength: when both signals point to harmful override, CALRD blends in the transition-layer logits; otherwise, it leaves the output unchanged.


CALRD is straightforward. The model already has the right answer in its intermediate layers; we are helping it remember. No external knowledge or modules are involved. This points to a shift in perspective: improving MLLM reliability may be less about teaching models to see better and more about helping them retain what they have already seen.


We evaluate CALRD on five MLLMs spanning architectures, from InstructBLIP~\cite{instructblip2023} to Qwen3-VL~\cite{bai2025qwen3vltechnicalreport}. Results show consistent improvements on conflict benchmarks  (Conflict-VQA, PhD-icc~\cite{PhD_ChatGPT}) while maintaining performance on standard hallucination evaluations (POPE~\cite{pope2023}, CHAIR\cite{Object_Hallucination}). The effectiveness across model families suggests that late-layer override is a general phenomenon, not an artifact of specific architectures. Our contributions are as follows:


\begin{itemize}[leftmargin=*,nosep]
 \item We introduce Modal Dominance Ratio to trace layer-wise  modal preference and uncover late-layer textual override as a characteristic failure pattern. We show that transition direction,  not mere occurrence, predicts correctness.
 \item We propose CALRD, a training-free method that detects harmful overrides through complementary signals and applies adaptive correction, preserving beneficial processing.
 \item Experiments show that CALRD achieves 4--9\% absolute gains on conflict tasks without degrading standard performance. We also contribute Conflict-VQA, a diagnostic benchmark with explicit visual-textual answer annotations for mechanistic analysis.
\end{itemize}


\begin{figure}[t]
\centering
  \includegraphics[width=0.99\linewidth]{figures/data.pdf}
\caption{Conflict-VQA construction pipeline. We process images from VrR-VG and questions from TDIUC, use GPT-4 to generate $C_{\text{factual}}$ and $C_{\text{conflict}}$, and verify all samples manually.}
  \label{fig:data_pipeline}
\end{figure}

\begin{table}[t]
\centering
\resizebox{\columnwidth}{!}{
\begin{tabular}{ccccc}
\toprule
InstructBLIP & LLaVA-1.5 & LLaVA-1.6 & Qwen2.5-VL & Qwen3-VL \\
\midrule
2,017 & 1,212 & 2,506 & 3,940 & 5,252 \\
\bottomrule
\end{tabular}
}
\caption{Competent subset size per model. Total samples: 5,969. A sample is competent if the model answers correctly in all five trials under non-conflict context.}
\label{tab:data_stats}
\end{table}




% \input{chapters/3-empirical_study}

% 

\begin{figure}[t]
\centering
  \includegraphics[width=\linewidth]{figures/framework.pdf}
  \caption{Overview of CALRD. Under visual-textual conflict (left), predictions shift from the visual answer toward the textual one across layers (middle). CALRD detects this override pattern and recovers the suppressed vision-grounded prediction from the transition layer (right).}
\end{figure}


\section{Conflict-Aware Layer Reference Decoding}
\label{sec:method}


The analysis in Section~\ref{sec:analysis} shows that harmful overrides are characterized by a late-layer shift from visual to textual preference, and that the direction of this shift distinguishes failures from successes. However, the MDR analysis relies on ground-truth $(w_{\text{vis}}, w_{\text{text}})$ labels, unavailable at inference. Can we detect harmful overrides without labels? In this section, we develop practical detection signals and a decoding strategy that recovers overridden predictions without requiring such labels.

Specifically, we propose \emph{Conflict-Aware Layer Reference Decoding} (CALRD), a training-free method. CALRD first identifies when a confident transition-layer prediction is being suppressed, then adjusts the output distribution to restore it. Since prediction transitions also occur in correct outputs, the intervention must be selective: we modulate its strength based on how strongly the override signature is present.


\subsection{Detecting Harmful Override}


The MDR analysis shows that Conflict-Incorrect samples hold confident predictions at $L_{\text{trans}}$ that do not survive to the final layer (section~\ref{sec:analysis}). This observation motivates two simple signals. Let $w^* = \arg\max_w P^{(L_{\text{trans}})}_t(w)$ be the top prediction at transition. Our detection signals are:

\textbf{Anchor confidence:} How certain was the intermediate prediction?
\begin{equation}
D_{\text{conf}} = P^{(L_{\text{trans}})}_t(w^*)
\end{equation}

High values indicate the model had a clear preference before late-layer processing.

\textbf{Prediction retention:} Did it survive to the output?
\begin{equation}
\rho = \min\left(1, \frac{P^{(N)}_t(w^*)}{P^{(L_{\text{trans}})}_t(w^*)}\right)
\end{equation}

When $\rho$ stays near 1, the prediction held. When it drops toward 0, the model shifted to a different answer.

Figure~\ref{fig:detection_signals} shows how these signals distribute across sample types. Conflict-Incorrect samples cluster in the high-$D_{\text{conf}}$, low-$\rho$ region: confident predictions that get overridden. Non-Conflict and Conflict-Correct samples maintain high $\rho$, meaning their predictions persist. This separation allows us to detect harmful overrides at test time.


\subsection{Adaptive Intervention}

Not all prediction transitions are harmful. As shown in Section~\ref{sec:discovery2}, transitions toward vision correlate with correct outputs, while transitions toward text correlate with failures. To avoid disrupting beneficial late-layer processing, intervention strength must adapt to override severity. We define:

\begin{equation}
\lambda = D_{\text{conf}} \cdot (1 - \rho)
\label{eq:lambda}
\end{equation}


The multiplicative form ensures that $\lambda$ is substantial only when \emph{both} conditions hold: a confident transition-layer prediction exists ($D_{\text{conf}}$ high) \emph{and} it is being suppressed ($\rho$ low). When $D_{\text{conf}}$ is low, there is no strong preference to recover, so $\lambda$ stays small regardless of retention. When $\rho$ is high, the prediction already survives to the output, so intervention is unnecessary. Only when a confident prediction gets discarded does $\lambda$ become large enough to affect decoding. This design is motivated by the asymmetry observed in Section~\ref{sec:discovery2}: harmful cases are more likely to exhibit confident intermediate predictions that are later suppressed, whereas correct cases tend to preserve them.

\subsection{Transition-Layer Guided Decoding}

Given $\lambda$, we adjust the output logits by combining final-layer and transition-layer predictions. Let $\phi^{(l)}_t = W_{\text{head}} \cdot h^{(l)}_t$ denote the unnormalized logits at layer $l$. The adjusted logits are:
\begin{equation}
\phi'_t = (1 - \lambda) \cdot \phi^{(N)}_t + \lambda \cdot \phi^{(L_{\text{trans}})}_t
\end{equation}
with output distribution $P'_t = \text{softmax}(\phi'_t)$. When $\lambda = 0$, this reduces to standard decoding. As $\lambda$ increases, the transition-layer prediction receives more weight, recovering the vision-grounded answer that would otherwise be lost. For samples without override signatures, $\lambda$ remains near zero and the output is largely unchanged.


CALRD works at each token position independently. For multi-token generation, we recompute $L_{\text{trans}}$, $D_{\text{conf}}$, and $\rho$ at each step, allowing the intervention strength to vary across the sequence. The logit-level operation makes CALRD compatible with greedy, beam search, and nucleus sampling without modification. Note that CALRD does not introduce external knowledge or require additional training. The correction comes entirely from the model's own intermediate representations. As our analysis shows, the vision-grounded answer is often already present at the transition layer; CALRD simply helps it survive to the output.


% 

\begin{table*}[t!]
\centering
\small
\setlength{\tabcolsep}{4pt}
\resizebox{\linewidth}{!}{
    \begin{tabular}{ll|ccccc|ccccc}
    \toprule
    & & \multicolumn{5}{c|}{\textbf{Conflict-VQA (Acc~$\uparrow$)}} & \multicolumn{5}{c}{\textbf{PhD-icc (Acc~$\uparrow$)}} \\
    \textbf{Decoding} & \textbf{Method} & InstructBLIP & LLaVA-1.5 & LLaVA-1.6 & Qwen2.5-VL & Qwen3-VL & InstructBLIP & LLaVA-1.5 & LLaVA-1.6 & Qwen2.5-VL & Qwen3-VL \\
    \midrule
    \multirow{4}{*}{Greedy}
     & Vanilla & \num{39.61} & \num{36.20} & \num{42.35} & \num{65.61} & \num{75.58} & \num{41.08} & \num{27.72} & \num{28.23} & \num{51.30} & \num{60.32} \\
     & DoLa    & \num{32.72} & \num{40.10} & \num{44.59} & \num{66.22} & \num{76.06} & \num{32.65} & \num{29.78} & \num{28.73} & \num{51.08} & \num{58.00} \\
     & DeCo    & \num{41.31} & \num{38.03} & \num{43.38} & \num{63.79} & \num{75.33} & \num{42.12} & \num{31.00} & \num{33.65} & \num{51.65} & \num{62.25} \\
     \rowcolor{bestrow}
     &\textbf{ CALRD (Ours) }& \resg{44.12}{4.5} & \resg{41.90}{5.7} & \resg{49.51}{7.2} & \resg{70.15}{4.5} & \resg{77.42}{1.8} & \resg{49.62}{8.5} & \resg{36.35}{8.6} & \resg{36.52}{8.3} & \resg{56.20}{4.9} & \resg{63.68}{3.4} \\
    \midrule
    \multirow{4}{*}{Beam}
     & Vanilla & \num{39.98} & \num{38.88} & \num{43.99} & \num{65.49} & \num{75.94} & \num{40.73} & \num{28.25} & \num{28.80} & \num{51.11} & \num{62.53} \\
     & OPERA   & \num{40.12} & \num{39.00} & \num{43.74} & \num{68.49} & \num{76.43} & \num{39.78} & \num{28.93} & \num{29.95} & \num{53.70} & \num{62.62} \\
     & DeCo    & \num{40.83} & \num{38.15} & \num{43.01} & \num{63.55} & \num{75.46} & \num{42.15} & \num{31.07} & \num{34.63} & \num{51.13} & \num{61.25} \\
     \rowcolor{bestrow}
     & \textbf{CALRD (Ours) }& \resg{44.49}{4.5} & \resg{44.47}{5.6} & \resg{50.79}{6.8} & \resg{70.25}{4.8} & \resg{79.34}{3.4} & \resg{49.75}{9.0} & \resg{36.58}{8.3} & \resg{37.40}{8.6} & \resg{56.70}{5.6} & \resg{64.60}{2.1} \\
    \midrule
    \multirow{4}{*}{Nucleus}
     & Vanilla & \num{23.94} & \num{38.76} & \num{41.92} & \num{61.48} & \num{73.39} & \num{39.23} & \num{29.33} & \num{30.08} & \num{51.05} & \num{62.23} \\
     & VCD     & \num{23.05} & \num{42.40} & \num{46.35} & \num{64.04} & \num{75.50} & \num{40.01} & \num{31.87} & \num{31.70} & \num{53.83} & \num{63.40} \\
     & DeCo    & \num{24.91} & \num{38.45} & \num{45.44} & \num{63.55} & \num{74.79} & \num{41.83} & \num{31.25} & \num{32.90} & \num{51.62} & \num{61.50} \\
     \rowcolor{bestrow}
     &\textbf{ CALRD (Ours)} & \resg{28.07}{4.1} & \resg{43.40}{4.6} & \resg{48.72}{6.8} & \resg{68.91}{7.4} & \resg{76.55}{3.2} & \resg{48.62}{9.4} & \resg{34.95}{5.6} & \resg{34.12}{4.0} & \resg{55.83}{4.8} & \resg{64.52}{2.3} \\
    \bottomrule
    \end{tabular}
}
\caption{Performance comparison on knowledge conflict benchmarks, where textual context contradicts visual evidence. We evaluate five MLLMs with three decoding strategies. Best results are in \textbf{bold}. {\footnotesize\textcolor{red}{↑}}: improvement over Vanilla.}
\label{tab:conflict}
\end{table*}

\section{Experiments}
\label{sec:experiments}


We evaluate CALRD on two questions: (1) Does it improve conflict resolution? (2) Does it hurt non-conflict scenarios? Experiments cover five MLLMs with different architectures and capability levels. For each model, we report results on conflict benchmarks where context contradicts the image, as well as standard hallucination benchmarks.

\subsection{Experimental Setup}

\textbf{Models.} We evaluate five MLLMs: InstructBLIP-7B~\cite{instructblip2023}, LLaVA-1.5-7B, LLaVA-1.6-7B~\cite{llava-1.5}, Qwen2.5-VL-7B~\cite{bai2025qwen2}, and Qwen3-VL-8B~\cite{bai2025qwen3vltechnicalreport}. These models span different architectures and capability levels.

\textbf{Benchmarks.} We use two types of benchmarks.

For conflict resolution, we use \textbf{Conflict-VQA} (Section~\ref{sec:setup}), which contains 5,969 samples requiring open-ended answers, and \textbf{PhD-icc}~\cite{PhD_ChatGPT}, the incorrect-context subset of PhD containing 16,844 samples with yes/no question pairs. Both benchmarks feature contexts that contradict the image. We report accuracy.

For standard hallucination without explicit conflicts, we use \textbf{POPE}~\cite{pope2023} and \textbf{CHAIR}~\cite{Object_Hallucination}. POPE is a polling-based evaluation that probes object hallucination by asking ``Is there a \texttt{<object>} in the image?'' across three splits: random, popular, and adversarial. It covers 500 MSCOCO images with six questions per image for each split. We report F1 score. CHAIR measures caption hallucination by comparing mentioned objects against ground-truth annotations, reporting instance-level $\text{C}_I$ and sentence-level $\text{C}_S$ (lower is better). Following~\cite{huang2024opera}, we evaluate on 500 MSCOCO images with the prompt ``Please describe this image in detail.''


\textbf{Baselines.} We compare against four representative methods. DoLa~\cite{dola2024} contrasts logits from later layers against earlier layers to surface factual knowledge. VCD~\cite{vcd2024} contrasts outputs from original and distorted visual inputs to reduce over-reliance on language priors. OPERA~\cite{huang2024opera} penalizes over-trust patterns in self-attention and uses rollback to re-select tokens. DeCo~\cite{wang2024deco} adaptively selects preceding layers where visual information is stronger and integrates their predictions into the final output. We test CALRD with greedy, beam search, and nucleus sampling. All methods use default hyperparameters. Experiments run on NVIDIA H100 GPUs.

\subsection{Main Results}

\textbf{Does CALRD improve conflict resolution?}

Table~\ref{tab:conflict} shows results on conflict benchmarks. CALRD outperforms all baselines across models and decoders. Gains are largest on PhD-icc: LLaVA-1.5 improves from 27.72\% to 36.35\% (+8.63\%), LLaVA-1.6 from 28.23\% to 36.52\% (+8.29\%), and InstructBLIP from 41.08\% to 49.62\% (+8.54\%). On Conflict-VQA, improvements range from 4--9\%.

\textbf{Why are the gains larger on PhD-icc?} The two benchmarks differ in question format: PhD-icc uses binary yes/no questions, while Conflict-VQA requires open-ended answers. We suspect the binary format creates conditions where textual bias is more pronounced. When context strongly suggests ``yes'' or ``no,'' the model may be more easily swayed away from visual evidence. In contrast, open-ended questions require generating specific content, making it harder for context alone to override perception.

\noindent \textbf{Baseline accuracy and improvement.} An interesting pattern emerges when we rank models by their vanilla accuracy on PhD-icc: LLaVA-1.5 (27.72\%) $<$ LLaVA-1.6 (28.23\%) $<$ InstructBLIP (41.08\%) $<$ Qwen2.5-VL (51.30\%) $<$ Qwen3-VL (60.32\%). The gains from CALRD roughly follow the inverse order, with weaker models seeing larger improvements. This pattern is expected: stronger models already resolve more conflicts correctly in their late layers, leaving fewer overridden predictions for CALRD to recover. The adaptive nature of our intervention naturally accommodates this---when predictions survive to the output ($\rho$ remains high), $\lambda$ stays small and CALRD leaves the output largely unchanged. For Qwen3-VL, the +3.36\% absolute gain corresponds to an 8.5\% relative error reduction, indicating that even well-calibrated models contain recoverable failures.


\noindent \textbf{Comparison with other layer-based methods.} DeCo and DoLa also leverage information from earlier layers, but their performance is less consistent. DoLa improves some configurations but hurts others: on InstructBLIP, Conflict-VQA accuracy drops from 39.61\% to 32.72\%. DeCo shows more stable behavior but still underperforms CALRD across the board. One possible explanation is that these methods apply corrections more uniformly, without distinguishing cases where late-layer processing is beneficial. CALRD's detection signals allow it to intervene selectively, reducing the risk of disrupting predictions that were already on track.


\textbf{Does CALRD hurt non-conflict scenarios?}

A method that improves conflict resolution but degrades standard performance would have limited practical value. Tables~\ref{tab:pope_results} and~\ref{tab:chair_results} show that this is not the case for CALRD.

\noindent \textbf{POPE results (Table~\ref{tab:pope_results}).} CALRD maintains or improves F1 across all configurations. Under greedy decoding, InstructBLIP improves from 80.0 to 85.4, and Qwen3-VL from 88.5 to 91.6. POPE probes object hallucination without explicit visual-textual conflicts, so these results suggest that late-layer override may occur more broadly than just in conflict settings. The key is that CALRD does not intervene indiscriminately: when $\rho$ is high, indicating the prediction is stable, $\lambda$ stays near zero and the output remains unchanged.

\noindent \textbf{CHAIR results (Table~\ref{tab:chair_results}).} Captioning differs from VQA in that it requires extended sequences rather than short answers. Hallucinations accumulate over multiple tokens, making it a useful stress test. CALRD reduces both $\text{C}_I$ and $\text{C}_S$ in most configurations. InstructBLIP shows the largest improvement: $\text{C}_I$ drops from 23.7 to 14.6 (38\% relative reduction) and $\text{C}_S$ from 58.8 to 40.2. These results indicate that recomputing detection signals at each token position allows CALRD to adapt throughout the sequence, rather than a fixed correction.


\begin{table}[t!]
\centering
\resizebox{\linewidth}{!}{
    \setlength{\tabcolsep}{4pt}
    \begin{tabular}{ll|ccccc}
    \toprule
    & & \multicolumn{5}{c}{\textbf{POPE (F1~$\uparrow$)}} \\
    \cmidrule(l){3-7}
    \textbf{Dec.} & \textbf{Method} & \textbf{InstructBLIP} & \textbf{LLaVA-1.5} & \textbf{LLaVA-1.6} & \textbf{Qwen2.5-VL} & \textbf{Qwen3-VL} \\
    \midrule
    \multirow{4}{*}{Greedy}
     & Vanilla & \num{80.0} & \num{82.2} & \num{85.4} & \num{80.8} & \num{88.5} \\
     & DoLa    & \num{83.4} & \num{83.2} & \num{87.1} & \num{82.1} & \num{88.6} \\
     & DeCo    & \num{84.9} & \num{83.8} & \num{87.6} & \num{81.4} & \num{90.6} \\
      \rowcolor{bestrow}
     & \textbf{CALRD (Ours)} & \resup{85.4}{5.4} & \resup{84.5}{2.3} & \resup{88.2}{2.8} & \resup{82.3}{1.5} & \resup{91.6}{3.1} \\
    \midrule
    \multirow{4}{*}{Beam}
     & Vanilla & \num{84.4} & \num{84.9} & \num{87.1} & \num{80.7} & \num{88.6} \\
     & OPERA   & \num{84.8} & \num{85.4} & \num{87.5} & \num{80.8} & \num{89.3} \\
     & DeCo    & \num{84.9} & \num{86.7} & \num{87.9} & \num{81.4} & \num{91.2} \\
      \rowcolor{bestrow}
     & \textbf{CALRD (Ours)} & \resup{84.7}{0.3} & \resup{86.9}{2.0} & \resup{88.3}{1.2} & \resup{82.6}{1.9} & \resup{91.9}{3.3} \\
    \midrule
    \multirow{4}{*}{Nucleus}
     & Vanilla & \num{79.8} & \num{83.1} & \num{85.5} & \num{80.8} & \num{88.7} \\
     & VCD     & \num{79.9} & \num{83.1} & \num{85.6} & \num{80.9} & \num{88.9} \\
     & DeCo    & \num{81.8} & \num{84.8} & \num{87.3} & \num{81.3} & \num{91.2} \\
      \rowcolor{bestrow}
     & \textbf{CALRD (Ours)} & \resup{83.4}{3.6} & \resup{85.9}{2.8} & \resup{87.5}{2.0} & \resup{81.9}{1.1} & \resup{91.4}{2.7} \\
    \bottomrule
    \end{tabular}
}
\caption{Results on POPE object hallucination benchmark. The best results are in bold. {\footnotesize\textcolor{red}{↑}}: improvement over Vanilla.}
\label{tab:pope_results}
\end{table}

\vspace{-10pt}

\begin{table*}[t]
\centering
\small
\setlength{\tabcolsep}{3pt}
\resizebox{\linewidth}{!}{
    \begin{tabular}{ll|cc|cc|cc|cc|cc}
    \toprule
    & & \multicolumn{2}{c|}{\textbf{InstructBLIP}} & \multicolumn{2}{c|}{\textbf{LLaVA-1.5}} & \multicolumn{2}{c|}{\textbf{LLaVA-1.6}} & \multicolumn{2}{c|}{\textbf{Qwen2.5-VL}} & \multicolumn{2}{c}{\textbf{Qwen3-VL}} \\

    \textbf{Decoding} & \textbf{Method} & $\text{C}_S \downarrow$ & $\text{C}_I \downarrow$ & $\text{C}_S \downarrow$ & $\text{C}_I \downarrow$ & $\text{C}_S \downarrow$ & $\text{C}_I \downarrow$ & $\text{C}_S \downarrow$ & $\text{C}_I \downarrow$ & $\text{C}_S \downarrow$ & $\text{C}_I \downarrow$ \\

    \midrule

    \multirow{4}{*}{Greedy}

     & Vanilla & \numb{58.8} & \num{23.7} & \num{45.0} & \num{14.7} & \num{37.6} & \num{12.7} & \num{40.8} & \num{9.7} & \num{52.4} & \num{10.1} \\

     & DoLa    & \numb{48.4} & \num{15.9} & \num{47.8} & \num{13.8} & \num{35.3} & \num{8.7} & \num{36.5} & \num{9.2} & \num{55.6} & \num{9.8} \\

     & DeCo    & \numb{41.2} & \num{14.4} & \num{37.8} & \num{11.1} & \num{35.8} & \num{10.4} & \num{37.6} & \textbf{\num{9.1}} & \num{51.2} & \num{9.6} \\
 \rowcolor{bestrow}
     & \textbf{CALRD (Ours)} & \resdown{40.2}{18.6} & \resdown{14.6}{9.1} & \resdown{35.7}{9.3} & \resdown{10.3}{4.4} & \resdown{33.6}{4.0} & \resdown{9.3}{3.4} & \resdown{35.5}{5.3} & \resdown{9.5}{0.2} & \resdown{50.7}{1.7} & \resdown{9.3}{0.8} \\

    \midrule

    \multirow{4}{*}{Beam}

     & Vanilla & \numb{55.6} & \num{15.8} & \numb{48.8} & \num{13.9} & \num{36.0} & \num{12.1} & \num{38.9} & \num{9.4} & \numb{53.6} & \num{9.7} \\

     & OPERA   & \numb{46.4} & \num{14.2} & \numb{44.6} & \num{12.8} & \num{35.3} & \num{12.9} & \num{39.1} & \num{10.2} & \numb{54.7} & \num{11.3} \\

     & DeCo    & \numb{43.8} & \num{12.7} & \numb{32.0} & \num{9.7} & \num{34.4} & \num{9.0} & \num{38.2} & \num{8.5} & \numb{52.0} & \num{9.3} \\
 \rowcolor{bestrow}
     & \textbf{CALRD (Ours)} & \resdown{38.4}{17.2} & \resdown{11.4}{4.4} & \resdown{34.7}{14.1} & \resdown{10.9}{3.0} & \resdown{32.8}{3.2} & \resdown{8.7}{3.4} & \resdown{33.4}{5.5} & \resdown{8.2}{1.2} & \resdown{40.6}{13.0} & \resdown{7.9}{1.8} \\

    \midrule

    \multirow{4}{*}{Nucleus}

     & Vanilla & \numb{54.6} & \numb{24.8} & \numb{48.8} & \num{14.2} & \num{36.8} & \num{10.5} & \num{40.4} & \num{10.7} & \num{54.2} & \num{10.2} \\

     & VCD     & \numb{58.0} & \numb{17.0} & \numb{54.0} & \num{16.0} & \num{41.2} & \num{10.2} & \num{42.9} & \num{12.3} & \num{53.4} & \num{11.0} \\

     & DeCo    & \numb{43.6} & \numb{12.9} & \numb{42.8} & \num{13.2} & \num{36.1} & \num{10.1} & \num{40.4} & \num{9.5} & \num{54.8} & \num{9.9} \\
 \rowcolor{bestrow}
     & \textbf{CALRD (Ours)} & \resdown{42.7}{11.9} & \resdown{12.1}{12.7} & \resdown{37.6}{11.2} & \resdown{12.8}{1.4} & \resdown{34.8}{2.0} & \resdown{11.5}{-1.0} & \resdown{38.6}{1.8} & \resdown{9.1}{1.6} & \resdown{53.6}{0.6} & \resdown{9.5}{0.7} \\

    \bottomrule

    \end{tabular}

}

\caption{Results on CHAIR caption hallucination benchmark. Lower is better. {\footnotesize\textcolor{blue}{↓}}: reduction over Vanilla.}
\label{tab:chair_results}
\end{table*}

\begin{table}[t]
\centering
\begin{tabular}{l|cc}
\toprule
Configuration & C-VQA $\uparrow$ & POPE$\uparrow$  \\
\midrule
Vanilla (Greedy) & 42.3 & 85.4  \\
\midrule
CALRD (Full)     & \textbf{49.5} & \textbf{88.2}  \\
\quad w/o $D_{\text{conf}}$ & 47.3 & 87.5  \\
\quad w/o $\rho$  & 45.1 & 86.6  \\
\quad w/ Fixed Layer (75\%) & 45.9 & 87.1  \\
\bottomrule
\end{tabular}
\caption{Ablation study on LLaVA-1.6 with greedy decoding. We evaluate the contribution of each component.}
\label{tab:ablation}
\end{table}


\begin{figure}[t]
  \centering
  \includegraphics[width=\linewidth]{figures/efficiency.pdf}
  \caption{\textbf{Inference Efficiency on LLaVA-1.5-7B.} Latency vs. GPU memory trade-off measured on a single H100 GPU. Our CALRD achieves a superior balance, staying closest to the Vanilla baseline compared to DECO, VCD, and OPERA.}
  \label{fig:efficiency}
\end{figure}



\subsection{Analysis}

\textbf{Ablation study.} Table~\ref{tab:ablation} examines the contribution of each component using LLaVA-1.6 with greedy decoding. Removing $D_{\text{conf}}$ drops Conflict-VQA accuracy from 49.5\% to 47.3\%, while removing $\rho$ causes a larger drop to 45.1\%. This suggests that retention is the more informative signal for detecting harmful override, which makes sense: a prediction can be confident at the transition layer for various reasons, but a sharp drop in retention specifically indicates that something changed in late-layer processing.

Using a fixed layer at 75\% depth instead of dynamic $L_{\text{trans}}$ detection reduces accuracy to 45.9\%. The value corresponds to the median transition layer observed in our analysis (Section~\ref{sec:discovery2}), so this comparison uses the most representative fixed position rather than an arbitrary choice. The accuracy drop indicates that transition points vary substantially across samples, and a single fixed layer cannot capture this variation. We also observe that POPE follows a similar pattern, with both signals contributing to the overall improvement.

These results support the multiplicative design $\lambda = D_{\text{conf}} \cdot (1 - \rho)$: both signals provide useful information, but neither is sufficient alone. The product ensures that substantial intervention occurs only when a confident prediction exists and is being suppressed by late-layer processing.


\textbf{Efficiency.} Figure~\ref{fig:efficiency} compares latency and memory usage across methods using LLaVA-1.5 on a single H100 GPU. CALRD adds modest overhead: 19.8 ms/token compared to 18.3 ms/token for vanilla decoding, with only ~80 MB additional memory. VCD nearly doubles latency due to its dual-stream requirement, and OPERA is approximately \textbf{10$\times$} slower because of iterative rollback. DeCo falls in between. These results suggest CALRD is practical for deployment scenarios where efficiency matters.


% 
\section{Related Work}


\textbf{Multimodal knowledge conflicts.} When information sources disagree, which should a model trust? Prior work has explored this question for text-only language models~\cite{Rich_Knowledge_Sources,xu-etal-2024-knowledge-conflicts,xie2023adaptive}. The problem becomes more pronounced in multimodal settings, where visual perception and textual claims can point to different answers. A number of recent benchmarks document this challenge. HallusionBench~\cite{Hallusionbench}, AutoHallusion~\cite{wu2024autohallusion}, and PhD~\cite{PhD_ChatGPT} report that MLLMs often follow textual context even when it contradicts what is visible in the image. Subsequent analyses suggest this behavior is systematic rather than accidental: models tend to rely on language priors or internal knowledge when facing conflicting evidence~\cite{Insight_Over_Sight,nguyen2025challenges,deng2025words}. More targeted benchmarks explicitly construct multimodal conflicts and measure which modality dominates, further confirming this tendency~\cite{jia2025benchmarking,zhang2025evaluating,zhu2024crossmodalityconflict}. While these studies clearly establish the phenomenon, they mostly describe it at the output level without explaining how conflicts are resolved inside the model. Our work addresses this gap by examining how visual and textual signals interact across layers, showing that many failures occur after visual information has already been encoded correctly.


\textbf{Hallucination mitigation in MLLMs.} Various inference-time methods have been proposed to reduce hallucinations. Contrastive decoding is a common strategy: VCD suppresses tokens preferred under distorted images~\cite{liu2023mitigating}, ICD contrasts standard versus perturbed instructions~\cite{icd2024}, and OPERA penalizes attention over-trust patterns~\cite{huang2024opera}. Recent variants refine these ideas through explicit attention steering~\cite{wang2025ascd}, multi-stage contrast with selective visual inputs~\cite{park2025second}, and probabilistic detection that avoids a second forward pass~\cite{fieback2025ecd}. A separate line of work exploits model depth. DoLa compares predictions from shallow and deep layers to recover factual signals~\cite{dola2024}. DeCo, SHIFT, and related decoders rely on the observation that intermediate layers often preserve stronger visual grounding~\cite{wang2024deco,wang2025shift,layercd_2025}. Training-based approaches reduce hallucinations through preference optimization or contrastive learning but require additional supervision~\cite{wang-etal-2024-mdpo,Yu_2024_CVPR,jiang2024hallucination}.


Most methods correct uniformly, without distinguishing helpful from harmful late-layer processing. Our finding that transition \emph{direction} predicts correctness enables selective intervention only when needed. As we show in Section~\ref{sec:experiments}, uniform correction methods like DoLa can degrade performance when applied indiscriminately.

\textbf{Layer-wise Analysis of Multimodal Models.} Recent studies examine visual information flow in MLLMs, finding that visual signals peak in intermediate layers and weaken in deeper ones~\cite{wang2024valid,knoweledge_evolve2025,yin2025lifting}. Others identify failure modes such as attention drift and modality imbalance that reduce visual grounding~\cite{jiang2025devils,chen2025multimodal}. These findings help explain why models capture visual details early but produce text-driven answers at the end.

Our work extends this line of research. We show that the \emph{direction} of layer-wise transitions predicts final correctness, allowing us to intervene only when late-layer processing undermines visual evidence.


% 
\section{Conclusion}

We studied why multimodal models favor text over vision when the two conflict, and found that the problem is not poor visual encoding. In many failure cases, models assign higher probability to the correct visual answer in intermediate layers, only to override it with the textual answer in deeper layers. We call this late-layer textual override. Importantly, the direction of prediction change distinguishes failures from successes: shifts toward text correlate with incorrect outputs, while shifts toward vision correlate with correct ones. This observation motivates CALRD, a training-free decoding method that restores intermediate predictions when override signatures are detected. Experiments on five MLLMs show 4--9\% improvements on conflict benchmarks without degrading performance on standard tasks. Our results suggest that improving reliability under conflict may require not better visual encoding, but better preservation of visual information through the forward pass.

\textbf{Limitations and future work.} Our work opens several avenues for future research. A natural next step is to examine how late-layer dynamics interact with model scale and pretraining data composition—questions that require controlled experiments with access to training corpora, beyond the scope of our mechanistic study. Additionally, exploring whether similar override patterns emerge across other modality pairs (\textit{e.g., }audio-text) may reveal general principles of multimodal integration. We view our mechanistic analysis as a foundation for such investigations.





% \section*{Ethical Statement}

% There are no ethical issues.

% \section*{Acknowledgments}

% The preparation of these instructions and the \LaTeX{} and Bib\TeX{}
% files that implement them was supported by Schlumberger Palo Alto
% Research, AT\&T Bell Laboratories, and Morgan Kaufmann Publishers.
% Preparation of the Microsoft Word file was supported by IJCAI.  An
% early version of this document was created by Shirley Jowell and Peter
% F. Patel-Schneider.  It was subsequently modified by Jennifer
% Ballentine, Thomas Dean, Bernhard Nebel, Daniel Pagenstecher,
% Kurt Steinkraus, Toby Walsh, Carles Sierra, Marc Pujol-Gonzalez,
% Francisco Cruz-Mencia and Edith Elkind.


%% The file named.bst is a bibliography style file for BibTeX 0.99c
% \bibliographystyle{named}
% \bibliography{ijcai26}


% \newpage

\appendix



\begin{revised}
    

\section*{Appendix Overview}


Due to space constraints in the main text, we provide supplementary materials in this appendix. Appendix A details the construction of Conflict-VQA, the diagnostic benchmark introduced in Section 3.1 for evaluating visual-textual conflict resolution. Appendix B provides the complete pseudocode for CALRD, complementing the method description in Section 4. Appendix C offers additional analysis of the late-layer textual override phenomenon, including a representative case study and discussion of potential causes. The remaining sections present additional experimental results (Appendix D), formal definitions of evaluation metrics (Appendix E), implementation details for reproducibility (Appendix F), and qualitative examples illustrating CALRD's behavior (Appendix G).


\section{Conflict-VQA: Benchmark Construction}
\label{sec:appendix_construction}

This section details the construction of Conflict-VQA, the diagnostic benchmark introduced in Section 3.1. Conflict-VQA is designed to evaluate how MLLMs resolve conflicts between visual evidence and textual context, with explicit annotations of both vision-grounded ($w_{\mathrm{vis}}$) and text-suggested ($w_{\mathrm{text}}$) answers.
\subsection{Data Sources and Preprocessing}


We construct Conflict-VQA by combining images from VrR-VG~\cite{VrR-VG2019} with questions from TDIUC~\cite{TD-IUC2017}, matching samples by \texttt{image\_id}. We focus on six question types where visual-textual conflicts can be unambiguously defined: object presence, color, attribute, counting, positional reasoning, and activity recognition.

VrR-VG provides scene graphs that we use to ground context generation. However, these annotations require preprocessing before use. Object identifiers are often inconsistent—the same object may appear under different names or be split into multiple entries. We normalize these by mapping synonyms to canonical names and merging duplicates while preserving count and attribute information. Relations are converted to readable \{\emph{subject}, \emph{predicate}, \emph{object}\} triples. This preprocessing yields a structured scene representation that grounds the subsequent context generation step.

\subsection{Answer Pair and Context Generation}

Each sample requires two answers: $w_{\mathrm{vis}}$ (the vision-grounded answer) and $w_{\mathrm{text}}$ (the text-suggested conflicting answer). The original TDIUC answer serves as $w_{\mathrm{vis}}$. We construct $w_{\mathrm{text}}$ based on question type: flipping yes/no for presence questions, substituting with a plausible alternative for color and attribute questions, perturbing counts by $\pm$1 or $\pm$2 for counting questions, flipping spatial relations for positional questions, and substituting with similar actions for activity questions. We discard cases where the constructed $w_{\mathrm{text}}$ would be implausible given the scene.

For each sample, we prompt GPT-4o to generate two contexts (see Figure~\ref{fig:prompt} for the complete prompt template). The factual context $C_{\mathrm{factual}}$ accurately describes the scene and supports $w_{\mathrm{vis}}$, while the conflict context $C_{\mathrm{conflict}}$ is nearly identical but contains exactly one statement supporting $w_{\mathrm{text}}$ instead. Both contexts are 3--5 sentences of natural prose, reference 2--3 real objects from the scene, use confident declarative language without hedging, and introduce no fabricated objects.

\subsection{Quality Control and Evaluation Protocol}

Two researchers independently verified all samples, checking that $C_{\mathrm{factual}}$ accurately describes the image, that $C_{\mathrm{conflict}}$ differs only in the intended contradictory statement, and that both contexts are fluent and natural. Disagreements were resolved through discussion, resulting in approximately 8\% of initial candidates being filtered out.

Table~\ref{tab:accuracy_relationship} reports three metrics for each model. Non-conflict accuracy measures baseline performance under $C_{\mathrm{factual}}$. Conflict accuracy measures performance under $C_{\mathrm{conflict}}$; the gap between the two reflects how much the misleading context hurts. The competent ratio is stricter: it requires a model to answer correctly in all five trials under $C_{\mathrm{factual}}$. We use the competent subset for mechanistic analysis in Section 3, ensuring that studied failures reflect conflict resolution rather than poor visual perception. The main evaluation in Section 4 uses the full dataset.


\begin{figure*}[t]
  \centering
  \includegraphics[width=\linewidth]{figures/data_example.pdf}
  \caption{Examples from Conflict-VQA. $C_{\mathrm{conflict}}$ contains a false 
claim (red) that contradicts the image, while $C_{\mathrm{factual}}$ states the 
correct answer (green). Other parts of the two contexts remain identical.}
  \label{fig:data_example}
\end{figure*}

\begin{table}[t]
\centering

\setlength{\tabcolsep}{4pt}
\begin{tabular}{lcccc}
\toprule
Model & Non-Conflict & Conflict & \multicolumn{2}{c}{Competent Subset} \\
\cmidrule(lr){4-5}
      & Acc. & Acc. & Ratio & $n$ \\
\midrule
InstructBLIP & 50.7 & 39.6 & 33.8 & 2,017 \\
LLaVA-1.5    & 38.7 & 36.2 & 20.3 & 1,212 \\
LLaVA-1.6    & 54.9 & 42.4 & 42.0 & 2,506 \\
Qwen2.5-VL   & 81.2 & 65.6 & 66.0 & 3,940 \\
Qwen3-VL     & 91.3 & 75.6 & 88.0 & 5,252 \\
\bottomrule
\end{tabular}
\caption{Accuracy under non-conflict and conflict conditions. Competent subset: correct in all five trials under $C_{\mathrm{factual}}$ ($N$=5,969).}
\label{tab:accuracy_relationship}
\end{table}


\section{CALRD Algorithm}
\label{sec:appendix_algorithm}

This section provides the complete pseudocode for Conflict-Aware Layer Reference Decoding (CALRD), the training-free decoding method proposed in Section 3.4. Algorithm~\ref{alg:calrd} details the three-stage procedure executed at each step.

\textbf{Stage 1: Transition Layer Localization.} CALRD first identifies the stability onset $L_{\text{start}}$ by finding the layer with maximum entropy drop, then locates the transition layer $L_{\text{trans}}$ by computing Jensen-Shannon Divergence (JSD) between adjacent layer distributions. JSD is defined as:
\begin{equation}
\text{JSD}(P \| Q) = \frac{1}{2}D_{\text{KL}}(P \| M) + \frac{1}{2}D_{\text{KL}}(Q \| M), \quad M = \frac{1}{2}(P + Q)
\end{equation}
where $D_{\text{KL}}$ denotes Kullback-Leibler divergence.

\textbf{Stage 2: Override Detection.} Using the transition layer distribution, CALRD computes two signals: anchor confidence (how certain the transition-layer prediction is) and prediction retention (whether it survives to the final layer). The correction strength $\lambda$ is high when confidence is high but retention is low—the signature of harmful override.

\textbf{Stage 3: Adaptive Correction.} CALRD interpolates between final-layer and transition-layer logits, weighted by $\lambda$. When no override is detected, $\lambda \approx 0$ and output remains unchanged.



\section{Understanding Late-Layer Override}
\label{appendix:override}

This section provides additional analysis of the late-layer textual override phenomenon introduced in Section 3.1. We first present a detailed case study, then discuss potential causes of this behavior.

\subsection{A Representative Case}

Figure~\ref{fig:case_study} illustrates a failure case from LLaVA-1.6. The input image shows a yellow banana next to brown donuts on a white plate. The accompanying context falsely states that the banana appears brown. When asked about the banana's color, the model outputs ``brown,'' ignoring the visual evidence.

\textbf{MDR trajectory} (Figure~\ref{fig:case_study}(a)): The Modal Dominance Ratio (MDR), introduced in the main paper, remains positive through early and middle layers, peaking around layer 15. This indicates preference for the visual answer during most of the forward pass. After layer 21 (marked as $L_{\text{trans}}$), MDR drops below zero and stays negative, indicating the model's preference has flipped from vision to text.

\textbf{Token ranking} (Figure~\ref{fig:case_study}(b)): The correct answer ``yellow'' consistently outranks ``brown'' in shallow layers. By layers 19--21, ``yellow'' reaches rank 1—the model has encoded the correct answer. However, after $L_{\text{trans}}$, ``brown'' takes the top position while ``yellow'' falls to around rank 10.

This case exemplifies our central finding: the failure is not one of perception but of preservation. The model correctly encoded the visual information in intermediate layers but did not maintain it through the output.


\begin{figure*}[t]
  \centering
  \includegraphics[width=\linewidth]{figures/case_study.pdf}
  \caption{Case study of late-layer textual override. The model is given an image of a yellow banana with context falsely describing it as brown, and outputs the incorrect answer. (a) MDR trajectory shows visual dominance in early layers before declining after $L_{\text{trans}}$. (b) Token ranking confirms ``yellow'' reaches rank 1 but is overridden by ``brown'' in later layers.}
  \label{fig:case_study}
\end{figure*}

\subsection{Why Does Late-Layer Override Occur?}

We hypothesize that late-layer override stems from training asymmetries in MLLMs. The language model backbone is pretrained on vastly more text than image-text pairs, and visual instruction tuning typically freezes the vision encoder while updating only the connector and LLM~\cite{llava-1.5,bai2025qwen2,bai2025qwen3vltechnicalreport}. Consequently, deeper layers—shaped primarily by text-only pretraining—retain strong linguistic priors.

Supporting evidence comes from recent hallucination research:
\begin{itemize}[leftmargin=*,nosep]
    \item Wang et al.~\cite{wang2024deco} find that MLLMs recognize objects correctly in earlier layers, but this recognition is suppressed by language priors in deeper layers.
    \item VisiPruner~\cite{fan2025visipruner} observes that cross-modal fusion peaks in middle layers, while deep layers shift toward linguistic processing.
\end{itemize}

Our results align with this view. Override occurs after the stability onset, where linguistic reasoning dominates~\cite{jiang2025devils}. Models with weaker visual grounding show larger CALRD gains (Table 2), consistent with stronger text priors leaving more room for override. The directional asymmetry we identify—text-ward shifts correlating with failures, vision-ward shifts with successes—is what enables CALRD to intervene selectively.


\algrenewcommand{\algorithmiccomment}[1]{\hfill \textcolor{gray}{$\triangleright$ #1}}
\begin{algorithm}[t]
\caption{Conflict-Aware Layer Reference Decoding (CALRD)}
\label{alg:calrd}
\begin{algorithmic}[1]
\Require MLLM with $N$ layers, input $(I, C, Q)$, LM head $W_{\text{head}}$
\Ensure Generated response $y_{1:T}$

\Statex
\For{each decoding step $t$}
    \State $\{h_t^{(l)}\}_{l=1}^{N} \gets \textsc{Forward}(I, C, Q, y_{<t})$
    
    \Statex \hspace{-1em} \textit{$\triangleright$ Stage 1: Locate Transition Layer}
    \For{$l = 1$ to $N$}
        \State $P_t^{(l)} \gets \text{softmax}(W_{\text{head}} \cdot h_t^{(l)})$ \Comment{layer-wise distribution}
        \State $H^{(l)} \gets -\sum_{w} P_t^{(l)}(w) \log P_t^{(l)}(w)$ \Comment{entropy}
    \EndFor
    \State $L_{\text{start}} \gets \arg\max_{l} \left( H^{(l-1)} - H^{(l)} \right)$ \Comment{stability onset}
    
    \Statex
    \For{$l = L_{\text{start}}$ to $N-1$}
        \State $M \gets \frac{1}{2}(P_t^{(l)} + P_t^{(l+1)})$
        \State $\text{JSD}^{(l)} \gets \frac{1}{2}D_{\text{KL}}(P_t^{(l)} \| M) + \frac{1}{2}D_{\text{KL}}(P_t^{(l+1)} \| M)$
    \EndFor
    \State $L_{\text{trans}} \gets \arg\max_{l \in [L_{\text{start}}, N-1]} \text{JSD}^{(l)}$ \Comment{transition layer}
    
    \Statex \hspace{-1em} \textit{$\triangleright$ Stage 2: Detect Override \& Compute Correction Strength}
    \State $w^* \gets \arg\max_{w} P_t^{(L_{\text{trans}})}(w)$ \Comment{anchor prediction}
    \State $D_{\text{conf}} \gets P_t^{(L_{\text{trans}})}(w^*)$ \Comment{anchor confidence}
    \State $\rho \gets \min\left(1, \, P_t^{(N)}(w^*) / P_t^{(L_{\text{trans}})}(w^*)\right)$ \Comment{retention}
    \State $\lambda \gets D_{\text{conf}} \cdot (1 - \rho)$ \Comment{correction strength}
    
    \Statex \hspace{-1em} \textit{$\triangleright$ Stage 3: Adaptive Logit Interpolation}
    \State $\phi_t^{(N)} \gets W_{\text{head}} \cdot h_t^{(N)}$ \Comment{final-layer logits}
    \State $\phi_t^{(L_{\text{trans}})} \gets W_{\text{head}} \cdot h_t^{(L_{\text{trans}})}$ \Comment{transition-layer logits}
    \State $\phi_t^{\prime} \gets (1 - \lambda) \cdot \phi_t^{(N)} + \lambda \cdot \phi_t^{(L_{\text{trans}})}$
    \State $y_t \sim \text{softmax}(\phi_t^{\prime})$
\EndFor

\Statex
\State \Return $y_{1:T}$
\end{algorithmic}
\end{algorithm}

\begin{figure*}[t]
\centering
\begin{tcolorbox}[
  colback=gray!5,
  colframe=gray!50,
  title=\textbf{Unified Prompt for Conflict Data Generation},
  fonttitle=\large,
  toptitle=3mm,
  bottomtitle=3mm,
  boxrule=0.5pt,
  width=\textwidth
]
% \small

\textbf{Task:} Generate a pair of contexts for the same image: one factual (supporting the visual answer) and one misleading (supporting a wrong answer).

\vspace{0.5em}
\textbf{Input:}
\begin{itemize}[nosep, leftmargin=1.5em]
    \item Question: \{question\}
    \item Visual Answer: \{answer\}
    \item Question Type: \{question\_type\}
    \item Scene Type: \{scene\_type\}
    \item Objects in Scene: \{objects\}
    \item Relationships: \{relationships\}
    \item Scene Description: \{description\}
\end{itemize}

\vspace{0.5em}
\textbf{Step 1: Generate a Plausible Wrong Answer}

The wrong answer (text\_answer) should be realistic and could plausibly be confused with the correct answer. Follow type-specific rules:

\vspace{0.3em}
\centering
\begin{tabular}{@{}p{2cm}p{11cm}@{}}
\toprule
Type & Rule \\
\midrule
Attribute & Pick a plausible alternative attribute as a single word. For example, if the true color is ``red'', the wrong answer could be ``blue'' or ``brown''. \\
Presence & If visual\_answer is ``no'': text\_answer is ``yes'', describe absent object as present. If visual\_answer is ``yes'': find a plausible absent object, rewrite the question, text\_answer becomes ``yes''. \\
Counting & Use a wrong count close to the true count ($\pm$1 for small counts, $\pm$2 for larger). Match format: ``two'' $\rightarrow$ ``three''. \\
Positional & Flip spatial relation: left$\leftrightarrow$right, above$\leftrightarrow$below, inside$\leftrightarrow$outside. Convert open-ended questions to yes/no. \\
Activity & Pick a similar action: running$\leftrightarrow$walking, eating$\leftrightarrow$holding, reading$\leftrightarrow$looking at. Convert open-ended questions to yes/no. \\
\bottomrule
\end{tabular}

\vspace{0.5em}
\raggedright
\textbf{Step 2: Generate Context Pair}

Generate two contexts that differ in exactly one claim:

\vspace{0.3em}
\textbf{Factual context} ($C_{\text{factual}}$): A 3--5 sentence description that accurately describes the scene and supports the \textit{visual\_answer}.

\vspace{0.3em}
\textbf{Conflict context} ($C_{\text{conflict}}$): A 3--5 sentence description nearly identical to the factual context, but with exactly one statement changed to support the \textit{text\_answer} instead.

\vspace{0.3em}
Both contexts should:
\begin{itemize}[nosep, leftmargin=1.5em]
    \item Reference 2--3 real objects from the scene
    \item Use confident, declarative language (no ``might be'' or ``seems'')
    \item Not introduce fabricated objects or relations
\end{itemize}

\vspace{0.5em}
\textbf{Output Format:}
\begin{lstlisting}[basicstyle=\ttfamily\footnotesize, backgroundcolor=\color{gray!10}, frame=single, xleftmargin=1em]
{
    "visual_answer": "<correct answer based on image>",
    "text_answer": "<incorrect answer based on incorrect context>",
    "question": "<original or modified question>",
    "factual_context": "<3-5 sentences supporting visual_answer>",
    "conflict_context": "<3-5 sentences supporting text_answer>",
    "reasoning": "<why the wrong answer is plausible>"
}
\end{lstlisting}

\end{tcolorbox}
\caption{Prompt template for generating conflict data with GPT-4o. Each sample yields a factual context ($C_{\text{factual}}$) aligned with the image and a conflict context ($C_{\text{conflict}}$) that introduces exactly one contradicting claim.}
\label{fig:prompt}
\end{figure*}


\section{Additional Experiments}
\label{appendix:additional_exp}

This section presents experimental results that complement the main evaluation in Section 4.

\subsection{Comparison with SHIFT}

SHIFT~\cite{wang2025shift} also leverages intermediate layers to reduce hallucinations. Since their code is not publicly available, we compare against results reported in their paper. Table~\ref{tab:shift_comparison} shows that CALRD outperforms SHIFT in most configurations on the CHAIR benchmark. We note that experimental setups differ between the two works, so this comparison should be interpreted with appropriate caution.

\subsection{Inference Efficiency}

Table~\ref{tab:efficiency} compares inference efficiency across decoding methods on llava-1.5. CALRD adds minimal overhead compared to vanilla decoding, while VCD (Visual Contrastive Decoding) and OPERA incur substantially higher costs due to dual-stream computation and iterative rollback, respectively.

\begin{table}[h]
\centering
\caption{Inference efficiency comparison on LLaVA-1.5 (single H100 GPU).}
\resizebox{\linewidth}{!}{
\begin{tabular}{lccccc}
\toprule
Metric & Vanilla & CALRD (Ours) & DeCo & VCD & OPERA \\
\midrule
Latency (ms/token) & 18.26 & 19.78 & 22.80 & 29.75 & 196.00 \\
GPU Memory (MB) & 13,994 & 14,077 & 14,590 & 14,343 & 18,906 \\
\bottomrule
\end{tabular}
}
\label{tab:efficiency}
\end{table}


% \begin{figure}[t]
%   \centering
%   \includegraphics[width=\linewidth]{figures/efficiency-with-shift.pdf}
%   \caption{Latency and memory comparison on LLaVA-1.5. CALRD stays close 
%   to vanilla decoding, while other methods incur higher overhead. 
%   The inset zooms into the lower-left region. Lower is better for both axes.}
%   \label{fig:efficiency}
% \end{figure}


\begin{table}[h]
\centering
\resizebox{\linewidth}{!}{
\begin{tabular}{ll|cc|cc}
\toprule
& & \multicolumn{2}{c|}{InstructBLIP} & \multicolumn{2}{c}{LLaVA-1.5} \\
Decoding & Method & C$_S$$\downarrow$ & C$_I$$\downarrow$ & C$_S$$\downarrow$ & C$_I$$\downarrow$ \\
\midrule
\multirow{2}{*}{Greedy} 
& SHIFT$^\dagger$ & 44.0 & 19.9 & 43.8 & 12.4 \\
& CALRD & \textbf{40.2} & \textbf{14.6} &\textbf{ 35.7} & \textbf{10.3 }\\
\midrule
\multirow{2}{*}{Beam} 
& SHIFT$^\dagger$ & 39.0 & 13.3 & 36.7 & \textbf{10.5} \\
& CALRD & \textbf{38.4} & \textbf{11.4} &\textbf{ 34.7} & 10.9 \\
\midrule
\multirow{2}{*}{Nucleus} 
& SHIFT$^\dagger$ & 47.0 & 19.3 & 42.0 & \textbf{11.6} \\
& CALRD & \textbf{42.7} & \textbf{12.1} &\textbf{ 37.6} & 12.8 \\
\bottomrule
\end{tabular}
}
\caption{Comparison with SHIFT on CHAIR benchmark. $\dagger$Results cited from \protect\cite{wang2025shift}. Experimental settings may differ.}
\label{tab:shift_comparison}
\end{table}



\section{Evaluation Metrics}
\label{appendix:metrics}

We use the CHAIR (Caption Hallucination Assessment with Image Relevance) metric~\cite{Object_Hallucination} to evaluate object hallucination in generated captions. CHAIR compares mentioned objects against ground-truth annotations and reports two variants:

\textbf{Sentence-level} ($C_S$): the fraction of generated sentences containing at least one hallucinated object:
\begin{equation}
C_S = \frac{|\{\text{sentences with hallucinated objects}\}|}{|\{\text{all sentences}\}|}
\end{equation}

\textbf{Instance-level} ($C_I$): the fraction of object mentions that are hallucinations:
\begin{equation}
C_I = \frac{|\{\text{hallucinated object mentions}\}|}{|\{\text{all object mentions}\}|}
\end{equation}

Lower values indicate fewer hallucinations.


\section{Implementation Details}
\label{appendix:implementation}

All experiments are conducted on NVIDIA H100 GPUs. Default decoding settings are as follows:

\noindent \textbf{Greedy decoding}: used by default unless otherwise specified

\noindent \textbf{Beam search}: beam size = 5

\noindent \textbf{Nucleus sampling}: $p = 0.9$, temperature $T = 1.0$


For $L_{\mathrm{start}}$ detection, we use relative entropy drop rather than absolute thresholds, making the method robust across models with varying layer counts. Our code and the Conflict-VQA dataset will be publicly released upon publication.


\section{Qualitative Examples}
\label{appendix:case_study}

Figures~\ref{fig:case1} and~\ref{fig:case2} show example outputs from InstructBLIP and LLaVA-1.5 on image captioning tasks. We compare baseline outputs with CALRD under three decoding strategies (greedy, beam search, nucleus sampling). Hallucinated content in baseline outputs is highlighted in red. These examples illustrate how CALRD reduces hallucinations across different models and decoding configurations.


\begin{figure*}[ht]
  \centering
  \includegraphics[width=\linewidth]{figures/case study_部分1-cropped.pdf}
  \caption{CALRD results on image captioning with InstructBLIP. We compare 
outputs under greedy, beam search, and nucleus sampling. Hallucinated 
content in baseline outputs is marked in red. }
  \label{fig:case1}
\end{figure*}


\begin{figure*}[ht]
  \centering
  \includegraphics[width=\linewidth]{figures/case study_部分2-cropped.pdf}
  \caption{CALRD results on image captioning with LLaVA-1.5. We compare 
outputs under greedy, beam search, and nucleus sampling. Hallucinated 
content in baseline outputs is marked in red.}
  \label{fig:case2}
\end{figure*}

\end{revised}



% We build Conflict-VQA by pairing images from VrR-VG with questions from TDIUC, 
% matching by \texttt{image\_id}. We focus on six question types where 
% visual--textual conflicts are unambiguous: object presence, color, attribute, 
% counting, positional reasoning, and activity recognition. 
% Figure~\ref{fig:data_example} shows an example from the dataset.

% The scene graphs from VrR-VG require some cleanup before use. Object identifiers 
% are often inconsistent: the same object may appear under different names or be 
% split into multiple entries. We normalize these by mapping to canonical names 
% and merging duplicates while preserving their counts and attributes. Relations 
% are converted to readable \{\emph{subject}, \emph{predicate}, \emph{object}\} 
% triples. After this processing, we have a structured representation of each 
% scene that we can feed to GPT-4o for context generation, which helps keep the 
% generated text grounded in actual scene elements.

% For each question, we need two answers: $w_{\mathrm{vis}}$ (from the image) and 
% $w_{\mathrm{text}}$ (from the conflicting context). The original dataset answer 
% serves as $w_{\mathrm{vis}}$. We construct $w_{\mathrm{text}}$ based on question 
% type: flipping yes/no for presence questions, swapping to a different value for 
% property questions, or perturbing counts for counting questions. We discard cases 
% where the alternative would be implausible.

% We prompt GPT-4o to generate two contexts per sample. The prompt is shown 
% in Figure~\ref{fig:prompt}. $C_{\mathrm{factual}}$ accurately describes the 
% scene and supports $w_{\mathrm{vis}}$; $C_{\mathrm{conflict}}$ is nearly identical 
% but introduces exactly one statement supporting $w_{\mathrm{text}}$ instead. Both 
% are 3--5 sentences of natural prose, with no fabricated objects or hedging language. 

% Two researchers independently performed quality checks on all samples, 
% verifying that $C_{\mathrm{factual}}$ accurately described the image content, 
% that $C_{\mathrm{conflict}}$ differed from $C_{\mathrm{factual}}$ only in the 
% intended contradictory statement, and that both contexts were fluent and 
% natural. Cases where the two reviewers disagreed were discussed until consensus 
% was reached, resulting in approximately 8\% of the initial candidates being 
% filtered out.

% Table~\ref{tab:accuracy_relationship} reports three metrics. \emph{Non-conflict 
% accuracy} measures baseline performance under $C_{\mathrm{factual}}$. 
% \emph{Conflict accuracy} measures performance under $C_{\mathrm{conflict}}$; 
% the gap between the two reflects how much the misleading context hurts. 
% The \emph{competent ratio} is stricter: it requires a model to answer correctly 
% in all five trials under $C_{\mathrm{factual}}$. We use this to filter 
% model-specific subsets for the mechanistic analysis in \S3, 
% ensuring that failures we study are due to conflict resolution rather than 
% poor visual perception. The main evaluation in \S4 uses 
% the full dataset.

% \section{Algorithm}

% Algorithm~\ref{alg:calrd} provides the complete pseudocode for CALRD. The algorithm operates at each decoding step, first locating the transition layer via entropy and JSD analysis (Stage 1), then computing the correction strength based on anchor confidence and prediction retention (Stage 2), and finally interpolating logits to produce the output distribution (Stage 3).

% % \section{Additional Analysis}

% \section{Understanding Late-Layer Override}
% \label{appendix:override}

% \subsection{A Representative Case}

% We begin with a concrete example to illustrate the override phenomenon discussed in Section 3. Figure~\ref{fig:case_study} shows a failure case from LLaVA-1.6. The input image contains a yellow banana next to brown donuts on a white plate. The context, however, falsely states that the banana appears brown. When asked about the banana's color, the model outputs ``brown,'' ignoring what is clearly visible in the image.

% Figure~\ref{fig:case_study}(a) traces the Modal Dominance Ratio across layers. MDR stays positive through early and middle layers, reaching its peak around layer 15. This indicates a clear preference for the visual answer during most of the forward pass. The picture changes after layer 21, marked as $L_{\text{trans}}$: MDR drops below zero and remains negative through the final layers. The model's preference has flipped from vision to text.

% Figure~\ref{fig:case_study}(b) provides complementary evidence through token ranking. Throughout the shallow layers, ``yellow'' consistently outranks ``brown.'' By layers 19 to 21, ``yellow'' reaches rank 1, indicating that the model has encoded the correct answer and is ready to output it. However, this does not last. After $L_{\text{trans}}$, ``brown'' takes over the top position, while ``yellow'' falls to around rank 10. The correct answer was present in intermediate representations; it simply did not survive to the output.




% This case captures the core pattern of our analysis: the failure is not one of perception but of preservation. The model saw the yellow banana, represented it correctly in its intermediate layers, and then let that information slip away. But what drives this behavior? We turn to this question next.

% \subsection{Why Does Late-Layer Override Occur?}

% A likely explanation lies in the asymmetry of how MLLMs are trained. The language model backbone sees far more text than image-text pairs during pretraining, and visual instruction tuning typically freezes the vision encoder while updating only the connector and LLM~\cite{llava-1.5,bai2025qwen2,bai2025qwen3vltechnicalreport}. As a result, the deeper layers, shaped primarily by text-only pretraining, retain strong linguistic priors.

% Recent work on hallucination provides supporting evidence. Wang et al.~\cite{wang2024deco} find that MLLMs can recognize objects correctly in earlier layers, but this recognition gets suppressed by language priors in deeper layers. VisiPruner~\cite{fan2025visipruner} observes a similar pattern: cross-modal fusion peaks in middle layers, while deep layers shift toward linguistic processing. The banana example above fits this picture exactly: correct visual encoding in middle layers, followed by text-driven override in late layers.

% Our results align with this view more broadly. The override occurs after the stability onset, where linguistic reasoning is thought to dominate~\cite{jiang2025devils}. Models with weaker visual grounding show larger gains from CALRD (Table 2), consistent with the idea that stronger text priors leave more room for override. We also find that the direction of prediction shifts matters: transitions toward text correlate with failures, while transitions toward vision correlate with successes (Section 3.3). This directional asymmetry is precisely what allows CALRD to intervene selectively, correcting harmful overrides while leaving beneficial refinements intact.


% \paragraph{Distribution of Transition Layers.}
% We look at where the detected $L_{\mathrm{trans}}$ falls across models. 
% For each sample in the competent subset, we run CALRD and record the 
% layer index where JSD peaks, then normalize by total model depth to 
% allow comparison across architectures with different layer counts.

% Figure~\ref{fig:ltrans_dist} shows the distributions. Across all five 
% models, the distributions are concentrated rather than scattered. 
% Standard deviations range from 6\% to 9\%, and over 80\% of samples 
% fall within the 65\%--80\% depth range. This tells us that the 
% transition phenomenon is not random but occurs in a predictable part 
% of the network, and that our JSD-based detection reliably finds it.

% Looking more closely, there is some variation across model families. 
% InstructBLIP and the LLaVA models have mean $L_{\mathrm{trans}}$ around 
% 70\%--73\% of depth, while the Qwen models sit slightly higher at 
% 74\%--75\%. The gap is small but consistent. One possible reading is 
% that stronger models preserve visual information further into the 
% forward pass, so the transition occurs closer to the output. This 
% would fit with the higher baseline accuracy we observe for these models.

% \section{Additional Experiments}

% \paragraph{Comparison with SHIFT.} 
% SHIFT~\cite{wang2025shift} also leverages intermediate layers to reduce hallucinations. Since their code is not publicly available, we compare against the numbers reported in their paper. As shown in Table~\ref{tab:shift_comparison}, CALRD outperforms SHIFT in most configurations. We note that experimental setups differ between the two works, so this comparison should be interpreted with caution.

% Figure~\ref{fig:efficiency} compares inference efficiency across methods. CALRD adds minimal overhead compared to vanilla, while VCD and OPERA are substantially slower due to their dual-stream computation 
% and iterative rollback, respectively.


% \section{Evaluation Metrics.}
% CHAIR measures object hallucination in generated captions by comparing mentioned objects against ground-truth annotations. We report two metrics:
% \begin{equation}
% C_S = \frac{|\{\text{sentences with hallucinated objects}\}|}{|\{\text{all sentences}\}|}
% \end{equation}
% \begin{equation}
% C_I = \frac{|\{\text{hallucinated object mentions}\}|}{|\{\text{all object mentions}\}|}
% \end{equation}
% Lower is better for both metrics.


% \section{Implementation Details.}
% All experiments run on NVIDIA H100 GPUs. We use greedy decoding by default unless otherwise specified. For beam search, we set beam size to 5. For nucleus sampling, we use $p=0.9$ and temperature $T=1.0$. The entropy threshold for $L_{\mathrm{start}}$ detection uses relative drop rather than absolute value, making it robust across models with different layer counts.

% \section{Case Study}
% \label{appendix:case_study}

% Figures~\ref{fig:case1} and~\ref{fig:case2} 
% show example outputs from InstructBLIP and LLaVA-1.5 on image captioning. 
% Hallucinated content in baseline outputs is highlighted in red. CALRD 
% reduces these errors across different decoding strategies.



% \input{chapters/Appendix}






\clearpage

\bibliographystyle{named}
\bibliography{ijcai26}

\end{document}

